
\subsection{Global Climate Commitments and Critical Minerals}
\label{sec:global climate commitments}
The global transition toward a low-carbon economy has significantly increased the strategic importance of critical minerals such as copper. International climate objectives established under the Paris Agreement require a rapid deployment of renewable energy systems, electric vehicles, electricity grids, and energy storage technologies, all of which depend heavily on copper-intensive infrastructure. As a result, several governments and international organizations have developed policies aimed at securing the long-term supply of critical minerals while reducing the environmental impact associated with their extraction and processing.

In parallel, many countries have adopted net-zero emissions targets and industrial decarbonization strategies. These policies increasingly emphasize the need to reduce greenhouse gas emissions across mineral supply chains while ensuring sufficient availability of raw materials required for the energy transition.


\subsection{Chilean Climate and Energy Policies}

Chile has established ambitious climate objectives aimed at supporting the global energy transition while reducing greenhouse gas emissions from its energy-intensive industries. Through its updated Nationally Determined Contribution (NDC) and Climate Change Framework Law, Chile has committed to achieving carbon neutrality by 2050 \cite{CAT2020}. The country also plans to progressively phase out coal-fired electricity generation and accelerate the expansion of renewable energy production.

Chile occupies a particularly favorable position in this transition because of its exceptional renewable energy potential, especially solar energy in the Atacama Desert and wind energy in coastal regions \cite{SolarAtlas2019}. Renewable sources represented nearly 70\% of Chile’s electricity generation in 2024 \cite{EnergiaEstrategica2025}. These developments are particularly important for the mining sector because electricity generation strongly influences the indirect emissions associated with copper extraction and mineral processing.


\subsection{Sustainability Standards in the Mining Sector}

Several international organizations and industry initiatives promote sustainability and emissions reduction in the mining sector. The International Council on Mining and Metals (ICMM) establishes environmental and climate-related principles for mining companies, including commitments related to emissions reduction, responsible resource management, and transparency. Similarly, the International Copper Association (ICA) supports research and sustainability initiatives related to copper production and lifecycle management.

More recently, certification frameworks such as the Copper Mark have emerged to encourage compliance with environmental, social, and governance (ESG) standards across the copper value chain. These initiatives aim to improve transparency, strengthen environmental accountability, and promote more sustainable mining practices through independent verification processes.


\subsection{Decarbonization Challenges for Copper Mining}

Despite the rapid expansion of renewable electricity and increasing sustainability commitments, several studies emphasize that decarbonizing copper mining remains a major industrial challenge. Copper demand is expected to increase significantly under global electrification scenarios, while declining ore grades and increasing extraction complexity continue to raise energy requirements in mining operations.

Consequently, future emissions reductions in the copper sector will likely depend on a combination of renewable electricity integration, electrification of mining equipment, improvements in process efficiency, and technological innovation. These interacting technological and structural constraints motivate the need for decomposition approaches capable of separating the different drivers influencing CO$_2$ emissions in copper mining.

\subsection{Relation Between Copper Production Decarbonization and Sustainable Development Goals}

While the United Nations Sustainable Development Goals (SDGs) were not specifically designed for the mining sector, copper production and its decarbonization are closely connected to several of these global objectives. As copper is an essential material for renewable energy systems, electric vehicles, electricity grids, and energy storage technologies, the sector plays a major role in supporting the global energy transition.

The transition toward lower-emission copper production aligns particularly well with the following SDGs:

\begin{itemize}
    \item \textbf{SDG 7 -- Affordable and Clean Energy:} promotes the expansion of renewable electricity generation and electrification technologies, both of which rely heavily on copper-intensive infrastructure.

    \item \textbf{SDG 9 -- Industry, Innovation and Infrastructure:} encourages the development of more sustainable industrial systems through technological innovation, process efficiency improvements, and cleaner mining operations.

    \item \textbf{SDG 12 -- Responsible Consumption and Production:} supports more sustainable resource management practices, including emissions reduction, energy efficiency, and recycling within the copper value chain.

    \item \textbf{SDG 13 -- Climate Action:} calls for urgent measures to reduce greenhouse gas emissions and mitigate climate change, objectives that are directly related to the decarbonization of mining activities.
\end{itemize}